\documentclass[UTF8,a4paper,12pt]{ctexart}
\usepackage{ctex}
\usepackage{amsmath}
\usepackage{amssymb}
\usepackage{amsthm}
\usepackage[margin=2.5cm]{geometry}

\setCJKmainfont{SimSun}
\setCJKsansfont{Microsoft YaHei}

\newtheorem{theorem}{定理}[section]
\newtheorem{lemma}{引理}[section]
\newtheorem{definition}{定义}[section]

\parindent = 2em
\linespread{1.4}

\title{英文/中文混排论文标题\\{\normalsize An English Subtitle Optional}}
\author{作者甲 \and 作者乙}
\date{\today}

\begin{document}
\maketitle

\begin{abstract}
本文研究……问题。We propose a novel method that…… 实验表明，所提方法在……数据集上取得了显著提升。

\textbf{关键词：} 深度学习；自然语言处理；中文排版
\end{abstract}

\section{Introduction}

研究背景与动机……

\section{相关工作}

\subsection{国外研究}

Overleaf~\cite{overleaf2024} 和 TeXstudio~\cite{texstudio} 是常用的 LaTeX 工具。

\subsection{国内研究}

中文文献引用示例……

\section{方法}

\begin{definition}[凸函数]
  设 $f: \mathbb{R}^n \to \mathbb{R}$，若对任意 $\mathbf{x}, \mathbf{y}$ 和 $\lambda \in [0,1]$，有
  \[
    f(\lambda \mathbf{x} + (1-\lambda)\mathbf{y})
    \leq \lambda f(\mathbf{x}) + (1-\lambda) f(\mathbf{y}),
  \]
  则称 $f$ 为凸函数。
\end{definition}

\begin{theorem}[极小值定理]
  设 $f$ 是紧集 $S$ 上的连续凸函数，则 $f$ 在 $S$ 上达到最小值。
\end{theorem}

\begin{proof}
  由 $f$ 在紧集上的连续性，$f$ 有界且能达到上下确界。凸性保证了最小值点的存在性。证明完毕。
\end{proof}

\section{实验}

实验设置与结果对比……

\section{结论}

总结与展望……

\begin{thebibliography}{9}
\bibitem{overleaf2024}
Overleaf Documentation, 2024.

\bibitem{texstudio}
TeXstudio Team. TeXstudio: A LaTeX editor. 2024.
\end{thebibliography}

\end{document}
